\section{问题一的模型建立与求解}

% ── 普通问题结构：问题与选择 → 关键数学关系 → 求解结果 → 直接答案
% ── 若本题为核心问题，额外增加：结果解释与机制、适用边界
% ── 删除本注释后填写实际内容

\subsection{问题分析与模型选择}

% 说明：为什么选择这个模型？有哪些候选方案，为何排除其他？
% 不要只写"采用 X 方法"，要写"采用 X 而非 Y，因为……"

采用 Z-score 标准化处理：$x' = \dfrac{x - \mu}{\sigma}$。

本问采用[方法名]，而非[候选替代]，原因在于[选择理由]。

\subsection{关键数学关系}

% 写核心公式和推导步骤，而非逐步罗列所有计算。

多元线性回归基础模型：
\[
  y = \theta_0 + \theta_1 x_1 + \theta_2 x_2 + \dots + \theta_n x_n + \epsilon
\]

\subsection{求解结果}

随机森林的测试集 $R^2$ 达到 0.896，优于线性回归的 0.834。

\subsection{结果解释与机制}

% 必填：结果为什么呈现当前形态？哪个因素是主要驱动？
% 例：$R^2$ 提升主要来自随机森林能捕捉非线性交互，而非简单参数数量增加。
% 禁止只写"可见模型效果良好"。

[结果意义]：[数值]表明[含义]，其主要原因是[机制]。
移除[某特征/约束]后，[指标]下降[幅度]，证明[该因素]的贡献是实质性的。

\subsection{适用边界}

% 必填：哪些假设如果改变，结论将随之变化？
% 例：上述结论在样本量 > 500 且特征无严重多重共线性时成立。

本问结论在[条件]下成立；若[关键假设变化]，则[结论将如何变化]。

\subsection{直接答案}

% 清晰回答题面要求的问题，可用表格或列表。
